\documentclass[a4paper]{article}
\usepackage[margin=3cm]{geometry}
\usepackage{graphicx}

\usepackage{hyperref}
\hypersetup{
	colorlinks=true,
	linkcolor=blue,
	filecolor=magenta,
	urlcolor=cyan,
	hyperindex=true,
	linktocpage=true,%makes the page number be the link in the index
}
\usepackage{listings}
\lstset{
	basicstyle=\tiny\ttfamily,
	columns=fullflexible,
	keepspaces=true,
	breaklines=true,
	frame=single,
	showstringspaces=true,
	tabsize=4,
	showspaces=true,
	showtabs=true,
}

\renewcommand{\arraystretch}{1.2}

\title{Antiproton endofullerenes}

\author{Chaya2766}

\date{last updated \today}

\begin{document}
	\sffamily
	\hyphenchar\font=-1
	\sloppy
	
	\maketitle
	
	\textbf{
		\sloppy
		\hyphenchar\font=-1
		}
	
	\vfill
	
	\begin{abstract}
		Investigation of antimatter storage using buckminsterfullerenes (A.K.A. C60 or Buckyballs) for use in the Stareater Expanse worldbuilding project.
		
		\medskip
		
		\textbf{Section 1} discusses whether or not antiproton endofullerenes are actually possible, as well as the reasons for my decision to include them in the setting.
		As of writing, stable long-term storage of antiprotons inside fullerenes is purely theoretical with mixed evidence, but might be experimentally verified in the future.
		I referenced a source that discusses this in much greater detail and also outlines an experiment that could verify the plausibility of this technology, as well as proposes a specific time period in the near future when the experiment could be ran.
		
		\medskip
		
		In \textbf{section 2} I ran simple calculations to get the following properties of bulk antiproton endofullerenes: density, specific energy, energy density, molar mass.
		I present and explain these calculations, and then show a concise table collecting the results.
		
		\medskip
		
		\textbf{Section 3} presents my calculations regarding the energy efficiency of producing antiproton endofullerenes, including a totally arbitrary assumption of the success rate of the antiproton being captured and contained in the fullerene.
		
		\medskip
		
		\textbf{Section 4} presents my calculations regarding efficiency of the antiproton endofullerenes as a rocket fuel, which I end up setting at around 1 million seconds Isp after exploring 2 approaches to calculating it.
	\end{abstract}
	
	\vfill
	
	\begin{quote}
		\begin{center}
			\textbf{DISCLAIMER AND WARNING}
		\end{center}
		
		This is not a scientific paper or engineering advice.
		This document is only intended as creative world‑building and has not been peer‑reviewed or approved for any practical application.
		Use of these ideas in real‑world designs or other serious application is entirely at the reader’s discretion and risk.
	\end{quote}
	
	%\noindent \rule{\linewidth}{1pt}
	\vfill
	
	\tableofcontents
	
	\pagebreak
	
	\section{Evidence of plausibility and worldbuilding potential}
	
	The evidence for and against the possibility of long term storage of antiprotons inside buckminsterfullerenes is discussed in much greater detail in the following paper:
	
	\medskip
	
	\noindent "Molecular Antimatter Confinement Experiment (MACE): Proposal to investigate synthesis of $\overline{p}$@C60 antimatter endohedral fullerenes by ion implantation at ELENA-TELMAX", April 29, 2025
	
	link: \href{https://cds.cern.ch/record/2930975/files/SPSC-P-372.pdf}{https://cds.cern.ch/record/2930975/files/SPSC-P-372.pdf}
	
	\medskip
	
	As a short, non-exhaustive summary, here are some of the points brought up:
	
	\begin{itemize}
		\item Some models of the buckminsterfullerene show that a point charge placed in the middle of the molecule would experience no net force, and as such would not be held in the middle.
		
		\item Experiments using a flourine ion, which is suggested as a good stand-in for an antiproton, showed that there is a potential well in the middle of the molecule, and the flourine ion does in fact get held in the middle away from the carbon atoms rather than experiencing no net force.
		
		\item No observation of an antiproton placed inside a buckminsterfullerene had been made so far, as such given the mixed theoretical evidence and total lack of practical evidence, the viability of antiproton endofullerenes cannot be accepted nor discarded.
	\end{itemize}
	
	The paper also describes an experiment that could verify whether antiprotons embedded inside fullerenes are stable or not, and the nearest opportunity to conduct this experiment is a time window stretching from late 2026 to early 2028, during which a particular needed piece of equipment capable of producing large numbers of antiprotons will not be in use, and as such could be available for the experiment.
	
	\begin{center}
		\rule{0.9\linewidth}{1pt}
		
		\medskip
		
		\textbf{usefulness in the stareater expanse setting}
	\end{center}
	
	Antiproton endofullerenes could be a very useful plot device resulting from a combination of two of their properties:
	
	\begin{itemize}
		\item Their production is incredibly inefficient, and as such their use is likely to be extremely situational.
		This means that their inclusion in the setting is likely to have only marginal effects and will likely not come into play in most scenarios.
		
		\item They may serve as an incredibly energy dense fuel for spacecraft, allowing them to make journeys to distant star systems at relativistic speeds.
		At the time of writing there is no technology established in the setting that would allow interstellar spacecraft to reach high fractions of the speed of light.
		The natural conclusion of this are small relativistic colony vessels departing to distant stars, leaving behind the powerful civilization of the solar system with no plausible way to go back, and arriving at their destination with relatively little equipment to re-seed the civilization with.
	\end{itemize}
	
	The opportunity presented by this is relevant in collaborative writing, in particular where representatives of the sol civilization are to meet a less advanced alien species.
	Direct contact of most fictional civilizations with the Stareater Expanse sol civilization would be decisively one-sided just due to the vast the scale difference between typical "reasonably powerful" sci-fi civilization and the usual hard sci-fi depiction of a civilization nearing the kardashev-2 status.
	Removing the solar system from the playing field of the first contact scenario and putting in its place a single vessel or small collection of vessels, cut off from the overwhelming wealth and experience of the solar system, is likely to make the odds significantly more even, as well as justify some level of inexperienced behaviour from the visitors in a plausible manner.
	It can also help justify retaining a third millenium level of technological advancement while accessing locations thousands of years distant from sol, such as omega centauri or others, which would have normally taken millions of years or more to reach by the natural slow expansion of the human civilization with technological advancement continuing throughout that time.
	
	\medskip
	
	In the end the negligible impact of antiproton endofullerenes on most aspects of the stareater expanse setting combined with the writing opportunities they produce makes overlooking their not yet scientifically proven status worth it.
	
	\bigskip
	\pagebreak
	
	\section{Properties of antiproton endofullerenes}
	
	\begin{center}
		\textbf{mass density}
	\end{center}
	
    Plain empty fullerenes have a density of 1650kg/m$^3$.
    Adding an antiproton into the middle does not increase the volume of the fullerene but does slightly increase the mass.
    
    \medskip
    
    There are 60 atoms of carbon in a fullerene, with each carbon atom on average massing 12.01u : the fullerene as a whole then masses $60\cdot12.01u = 720.6u$.
    
    \medskip
    
    Adding one proton increases the mass by 1u, from 720.6u to 721.6u, so the density should change by a factor of $\frac{721.6u}{720.6u} = 1.001387732$.
    As such the density of bulk antiproton endofullerenes is $1650kg/m^3 \cdot 1.001387732 = 1652.289759 kg/m^3$.
    
    \begin{center}
    	\textbf{specific energy}
    \end{center}
    
    When the shell of the antiproton endofullerene is disrupted, the antiproton will annihilate most likely against one of the carbon atoms. This converts the whole mass of the antiproton, as well as the mass of the proton it annihilated against, as energy (of what kind?) - thereby releasing the mass-energy of 2u.
    
    \medskip
    
    1u = 1.66053907 E-27 kg.
    
    \medskip
    
    $$ 1.66053907 E-27 kg \cdot c^2 = 1.66053907 E-27 kg \cdot 8.987551787 E+16 m^2/s^2 = 1.49241809 E-10 J $$
    
    \medskip
    
    1.49241809 E-10 J is released per buckyball, and again each antiproton endofullerene masses 721.6u (which is 1.19824499 E-24 kg).
    
    \medskip
    
    $$ \frac{1.49241809 E-10 J}{1.19824499 E-24 kg} = 1.245503301 E+14 J/kg = 124.5503301 TJ/kg $$
    
    \begin{center}
    	\textbf{energy density}
    \end{center}
    
    Multiply the specific energy of 1.245503301 E+14 J/kg with the density of 1652.289759 kg/m$^3$.
    
    \medskip
    
    $$ (1.245503301 E+14 J/kg) \cdot (1652.289759 kg/m^3) = 2.057932349 E+17 J/m^3 = 205.7932349 PJ/m^3 $$
    
    \begin{center}
    	\textbf{molar mass}
    \end{center}
    
    Can be calculated by simply multiplying the mass of a single antiproton endofullerene (721.6u = 1.19824499 E-24 kg) by the avogadro constant (6.02214076 E+23 mol$^{-1}$):
    
    $$ (1.19824499 E-24 kg) \cdot (6.02214076 E+23 mol^{-1}) = 0.7216 kg/mol $$
    
    \begin{center}
	    \begin{tabular}{| c | c |}
	    	\hline
	    	\multicolumn{2}{|c|}{\textbf{summary properties of antiproton fullerenes}} \\\hline
	    	density & 1652.289759 kg/m$^3$ \\\hline
	    	specific energy & 124.5503301 TJ/kg \\\hline
	    	energy density & 205.7932349 PJ/m$^3$ \\\hline
	    	molar mass & 0.7216 kg/mol \\\hline
	    \end{tabular}
    \end{center}
    
    \pagebreak
    
    \section{antiproton endofullerene production efficiency}
    
    I'm relying on a document that unfortunately cannot be easily shared as it requires an account on the Orion's Arm forum, but here is a link anyway: 
    
    \href{https://www.orionsarm.com/forum/attachment.php?aid=2348}{https://www.orionsarm.com/forum/attachment.php?aid=2348}
    
    For searching in other sources, it is titled "Antimatter induced fusion and thermonuclear explosions*" and credited to Andre Gsponer and Jean-Pierre Hurni from "Independent Scientific Research Institute".
    Release date is "February 2, 2008" and visible on the title page is also "ISRI-86-06.4".
    
    \medskip
    
    In section 8 the document makes an assumption that the power efficiency of the accelerators would be 25\% and claims that efficiency of the antiproton collection could be as high as 5\%.
    In my case, as I am not only collecting the antiprotons but additionally attempting to implant them into C60 fullerenes, I am adding a third factor that is the success rate of implanting the aniproton into the fullerene, versus the expected failure through the antiproton contacting one of the carbon atoms and annihilating.
    I am arbitrarily setting the expected success rate of antiproton capture at 0.1\%.
    
    \medskip
    
    All together the power efficiency of the accelerators, the collectors, and the capture success rate yield a final efficiency as such:
    
    $$ 25\% \cdot 5\% \cdot 0.1\% = 0.25 \cdot 0.05 \cdot 0.001 = 0.0000125 $$
    
    This means that every joule poured into antiproton endofullerene production results in an end product containing a total of 0.0000125 joules of energy stored in the form of antiproton endofullerenes.
    
    \medskip
    
    With this efficiency a 1-gigawatt plant operating for 24 hours would produce 1.08GJ worth of antiproton endofullerenes, which would mass 8.6712mg.
    To produce one kilogram of the product, 9.964 Exajoules would have to be spent, which is about 5.77 days worth of humanity's average power usage in gregorian year 2020 (about 2E+13 watts).
    
    \medskip
    
    Correspondingly, producing 1000 tonnes of the product would require 9.964 Yottajoules which is about 31.574 years worth of an entire kardashev-1 level civilization (1E+16 watts) or about 996.4 seconds of the power output of a kardashev-1.6 civilization, which in the Stareater Expanse setting is projected to be the state of humanity in a time period somewhere between years 2900 and 3000.
    
    \medskip
    
    This seems to be a sweet spot such that the overwhelming inefficiency of the process prohibits it for most use cases, but determined organizations in the late third millenium would be able to mass together enough antimatter fuel to launch missions to distant milky way objects of interest.
    
    \pagebreak
    
    \section{antiproton endofullerenes as rocket fuel}
    
    The total mass of an antiproton endofullerene is 721.6u and the annihilated mass is 2u, leaving behind 719.6u of mass that did not get annihilated.
    
    \medskip
    
    I will take two approaches to determining the resulting exhaust velocity.
    
    \medskip
    
    \textbf{1. simple kinetic energy conversion:} I will assume that the mass-energy of the annihilated proton and antiproton is dumped into the remaining 719.6u of mass as kinetic energy.
    
    $$ 2u \cdot c^2 = 2.9848362 E-10 J ; 719.6u = 1.194923914 E-24 kg $$
    
    Kinetic energy formula: $E_k = \frac{m \cdot v^2}{2}$
    
    Rearranged to deduce the velocity: $$v = \sqrt{\frac{2 E_k}{m}}$$
    
    $$ v = \sqrt{\frac{2 \cdot 2.9848362 E-10 J}{1.194923914 E-24 kg}} = 22351420.19 m/s $$
    
    Specific impulse: $\frac{22351420.19 m/s}{9.80665 m/s^2} = 2279210.555 s \approx 2.28 E+6 s$
    
    \medskip
    
    \textbf{2. heating:} the mass-energy of the annihilated particles is dumped into the remaining particles as heat, and the exhaust escapes with the corresponding thermal velocity.
    As to not bring nuclear physics into this I have to assume that despite having annihilated one of the regular protons, the exhaust is still 100\% carbon atoms, this should not have a major effect on the result.
    
    \medskip
    
    Carbon's specific heat capacity is 709.1 J/(kg·K).
    The energy from annihilation would be spread evenly between all 60 carbon atoms, giving each of them $\frac{2.9848362 E-10 J}{60} = 4.974727 E-12 J$ of thermal energy.
    
    \medskip
    
    A carbon atom masses 12.01u = 1.99430742 E-26 kg.
    That mass of carbon being given 4.974727 E-12 J of thermal energy would give temperature as such:
    
    $$ \frac{4.974727 E-12 J}{1.99430742 E-26 kg \cdot 709.1 J/(kg·K)} = 351778799382 K $$
    
    \medskip
    
    The root mean square velocity of particles in a hot gas can be calculated using the following equation:
    
    $$ v = \sqrt{\frac{k_B \cdot T}{m}} $$
    
    where:
    
    \begin{itemize}
    	\item $T$ is the temperature of the gas
    	
    	\item $m$ is the mass of the particle
    	
    	\item $k_B$ is the boltzmann constant = 1.380649 E-23 J/K
    \end{itemize}
    
    For carbon with particle mass 1.99430742 E-26 kg heated to 351 778 799 382 K it gives:
    
    $$ \sqrt{\frac{1.380649 E-23 J/K \cdot 351778799382 K}{1.99430742 E-26 kg}} = 15605598.17 m/s $$
    
    Specific impulse: $\frac{15605598.17 m/s}{9.80665 m/s^2} = 1591328.147 s \approx 1.59 E+6 s$
    
    \medskip
    
    I am inclined to include some inefficiency, but not too much, perhaps bringing the effective exhaust velocity down to 10 000 000 m/s and specific impulse to a bit over 1 million seconds, which would represent about 45\% to 64\% efficiency.
    
    \pagebreak
    
    Two equations for use in this section, the newtonian and relativistic versions of the rocket equation.
    
    $$ \Delta v = v_e \cdot ln(\frac{M_f}{M_e}) $$
    
    $$ \Delta v = c \cdot tanh\left( \frac{v_e}{c} \cdot ln(\frac{M_f}{M_e}) \right) $$
    
    where:
    
    \begin{itemize}
    	\item $\Delta v$ is the rocket's available velocity budget
    	
    	\item $v_e$ is the exhaust velocity of the rocket
    	
    	\item $M_f$ is the mass of the rocket when fully fueled
    	
    	\item $M_e$ is the mass of the rocket when empty
    	
    	\item $c$ is the speed of light = 299 792 458 m/s
    \end{itemize}
    
    Given the 10 000 000 m/s exhaust velocity, a spacecraft that carries enough fuel to equal its own mass would be able to reach roughly 2.312\% of the speed of light:
    
    $$ 10Mm/s \cdot ln(2) = 2.31209\%c \qquad ; \qquad 1c \cdot tanh\left( \frac{10Mm/s}{c} \cdot ln(2) \right) = 2.31168\%c $$
    
    One whose mass when fueled is 20 times greater than when empty would be able to reach nearly 10\% of the speed of light:
    
    $$ 10Mm/s \cdot ln(20) = 9.993\%c \qquad ; \qquad 1c \cdot tanh\left( \frac{10Mm/s}{c} \cdot ln(20) \right) = 9.96\%c $$
    
    And a spacecraft whose functional parts make up only 1\% of the mass, with the remainder being fuel, would be able to reach over 15\% of the speed of light:
    
    $$ 10Mm/s \cdot ln(100) = 15.361\%c \qquad ; \qquad 1c \cdot tanh\left( \frac{10Mm/s}{c} \cdot ln(100) \right) = 15.241\%c $$
    
    Stretching the limits a bit, a spacecraft that is 99.9\% fuel and 0.1\% hull would be able to reach well past 20\% of the speed of light, which means coasting speed of 10\%c even if the velocity budget has to be split in half between the launch and slowdown at the other end:
    
    $$ 10Mm/s \cdot ln(1000) = 23.0418\%c \qquad ; \qquad 1c \cdot tanh\left( \frac{10Mm/s}{c} \cdot ln(1000) \right) = 22.6425\%c $$
    
    At such high mass ratios using a staged rocket would be expected, and discarding expensive equipment such as superconducting magnetic nozzles wouldn't be much more of a sacrifice than using antimatter fuel already is.
    Advanced materials such as graphene and carbon nanotubes, combined with sufficient wealth and will, could likely implement even the last of these options.
    
    \medskip
    
    Although theoretically it could go on arbitrarily far, eg. letting there be 0.01\% hull with the rest being fuel allows 29.79\%c $\Delta V$ and going down to 0.001\% would enable 36.62\%, but practical limits have to be drawn somewhere and 10\%c is already a lot.
    
    \medskip
    
    Launch via laser sail or some other beamed energy system, if deemed viable, could even further speed up the trip by removing the need to use onboard fuel for launch, and saving all of it for slowdown and in-transit maneouvering.
	
\end{document}